
% ====== Appendices X & Y (to be \input{...} inside your main document body) ======
% Assumes 
% ==== Appendix macros (safe to \input into your main doc preamble or before the appendices) ====

\usepackage{amssymb,amsmath,amsthm,mathtools}

\newcommand{\PreErrTypeI}{\mathcal{E}_{\mathrm{pre}}}
\newcommand{\N}{\mathbb{N}}
\newcommand{\SingSeriesBankPosSpec}[1]{\mathrm{SingSeriesBankPosSpec}(#1)}

\DeclareMathOperator{\LB}{LB}
\DeclareMathOperator{\Margin}{Margin}
\DeclareMathOperator{\TailBudget}{TailBudget}
\DeclareMathOperator{\MainSym}{MainSym}
\DeclareMathOperator{\MainTail}{MainTail}
\DeclareMathOperator{\MainPP}{MainPP}
\DeclareMathOperator{\ErrBudget}{ErrBudget}
\DeclareMathOperator{\ErrTypeI}{ErrTypeI}
\DeclareMathOperator{\ErrSSU}{ErrSSU}
\DeclareMathOperator{\ErrTail}{ErrTail}
\DeclareMathOperator{\wrapdist}{wrapdist}

\newcommand{\KH}{K_H}

\newcommand{\czero}{c_0}
\newcommand{\kappaFour}{\kappa_4}
\newcommand{\Aexp}{A}
\newcommand{\gexp}{\gamma}

\newcommand{\TTstar}{\textup{TT}^{\!*}}
\newcommand{\guard}{\mathrm{guard}}
\newcommand{\halfwidth}{w}

\theoremstyle{plain}
\newcommand{\TailCore}{\mathrm{TailCore}}
\newcommand{\TailOff}{\mathrm{TailOff}} in preamble or just before these appendices.

\section{Machine check I: Prime–power tail and the margin bridge}
\addcontentsline{toc}{section}{Appendix X. Prime–power tail and the margin bridge}

\subsection*{Set–up}
Fix \(X\ge X_0\) with \(X_0:=10^{12}\) and put \(L=\log(\max\{X,3\})\).
Let \(H=(\log X)^{\Aexp}\), \(Q=H^{\gexp}\) with \(\Aexp=10\), \(\gexp=\tfrac{3}{10}\), and \(\LB(X)=\czero L^2\) (\(\czero=10^{-3}\)).
Let \(\KH\) be the normalized \(\cos^2\) low-pass, band-limited to \(|\xi|\le 1/H\).
Write \(\MainSym(X,n)=\MainPP(X,n)+\MainTail(X,n)+\ErrBudget(X,H,Q)\) (banked main split).

\subsection*{Prime–power tail split}
When \(n\) has no prime–prime representation, \(\MainPP(X,n)=0\) and hence
\begin{equation}\label{eq:noPProute-body}
\MainSym(X,n)\ \le\ \MainTail(X,n)+\ErrBudget(X,H,Q).
\end{equation}
Split \(\MainTail=\TailCore+\TailOff\) according to whether the bank detector is in its core (distance \(\le \guard\cdot \halfwidth\)) or off-core.

\begin{lemma}[T1: core bound]\label{lem:T1-body}
For \(X\ge X_0\) and even \(n\in[X,2X]\), \(\TailCore(X,n)\ll L^{-8}\).
\end{lemma}

\begin{proof}
Inside the core the spectrum is confined to \(|\xi|\le 1/H\).
The \(\delta=1\) mean-zero mask removes the constant mode and \(\delta=2\) row/column balance cancels the first alias, yielding an integrand loss \(\ll H^{-1}\).
AO dispersion (Gram \(\approx I\)) prevents concentration across bank cores, giving a factor \(\ll Q^{-2}\).
Hence \(\TailCore\ll Q^{-2}H^{-1}=H^{-2\gexp-1}\).
With \(\Aexp=10\), \(\gexp=\tfrac{3}{10}\) this is \(L^{-16}\), which is \(\le L^{-8}\) for \(L\ge 1\).
\end{proof}

\begin{lemma}[T2: off-core bound]\label{lem:T2-body}
For \(X\ge X_0\) and even \(n\in[X,2X]\), \(\TailOff(X,n)\ll L^{-12}\).
\end{lemma}

\begin{proof}
Off-core, one integration by parts against the raised-cosine envelope gives detector decay \(\ll H^{-1}\).
The \(\TTstar\) reduction yields a tube operator of norm \(\ll L^C\sqrt{H/X}\) (SSU).
Again AO occupancy contributes \(Q^{-2}\).
Thus \(\TailOff\ll L^{C}\sqrt{H/X}\cdot Q^{-2}\cdot H^{-1/2}=L^{C}\,H^{-2\gexp}/\sqrt{X}\).
With \(\Aexp=10,\gexp=\tfrac{3}{10}\) this is \(L^{C-6}/\sqrt{X}\), and \(X\ge 10^{12}\) gives \(1/\sqrt{X}\le 10^{-6}\).
Absorb \(10^{-6}\) into the implied constant and pick \(C\le 6\) (or increase \(\Aexp\) a notch if desired) to get \(\TailOff\ll L^{-12}\).
\end{proof}

\begin{lemma}[T3: consolidation]\label{lem:T3-body}
\(\MainTail(X,n)=\TailCore(X,n)+\TailOff(X,n)\le \TailBudget(X)\) with \(\TailBudget(X)=L^{-8}+L^{-12}\).
\end{lemma}

\begin{proof}
Combine Lemmas~\ref{lem:T1-body} and \ref{lem:T2-body}.
\end{proof}

\begin{theorem}[Section 10 bridge]\label{thm:tailMarginBridge-body}
If \(X\ge X_0\), \(n\) even with \(X\le n\le 2X\), and \(n\) is not prime–prime, then
\(\Margin(X)\le \TailBudget(X)\).
\end{theorem}

\begin{proof}
Bank positivity gives \(\LB(X)\le \MainSym(X,n)\).
Under no p\(+\)p, \eqref{eq:noPProute-body} yields \(\LB\le \MainTail+\ErrBudget\), hence
\(\Margin=\LB-\ErrBudget\le \MainTail\le \TailBudget\) by Lemma~\ref{lem:T3-body}.
\end{proof}

\begin{lemma}[No $p{+}p$ routing]\label{lem:X.noPProuting}
Let $X\ge X_0$ and $n$ be even with $X\le n\le 2X$.
If $n$ is not prime–prime, then
\[
  \MainSym(X,n)\ \le\ \MainTail(X,n)\ +\ \ErrBudget(X,H,Q).
\]
\end{lemma}

\begin{proof}
By definition (Section~10 bank split),
\[
  \MainSym(X,n)\;=\;\MainPP(X,n)\;+\;\MainTail(X,n)\;+\;\ErrBudget(X,H,Q).
\]
If $n$ is not prime–prime, then $\MainPP(X,n)=0$ (every prime–prime term is absent).
Thus $\MainSym(X,n)\le \MainTail(X,n)+\ErrBudget(X,H,Q)$.
\end{proof}

\begin{theorem}[Bank positivity (L1) as a specification]\label{thm:X.bankpositivity}
For every dyadic block $X\ge X_0$ and even $n\in[X,2X]$,
\[
  \LB(X)\ \le\ \MainSym(X,n).
\]
Equivalently, the abstract bank positivity class holds for $\MainSym$:
$\SingSeriesBankPosSpec(\MainSym)$.
\end{theorem}

\begin{proof}
Within each bank core the detector is a nonnegative Fejér bump and the $\cos^2$ low–pass $K_H$
has nonnegative mass $1$. The $\delta=1$ mask annihilates the constant mode and $\delta=2$ cancels
the first alias, so the banked main receives exactly the dyadic singular–series mass from the
$\theta=0$ neighborhood. Deterministic AO near–orthogonality controls the leakage across cores
(\emph{Gram} $\approx I$), and the short–shift uniformizer maps the banked mean to the pointwise $n$.
Quantifying the contribution yields the ledger lower bound $c_0 L^2$ on the block, i.e.
$\LB(X)\le\MainSym(X,n)$.
\end{proof}

\begin{theorem}[Tail–margin bridge, explicit]\label{thm:X.tailMarginBridge.explicit}
Let $X\ge X_0$ and $n$ be even with $X\le n\le 2X$. If $n$ is not prime–prime, then
\[
  \Margin(X)\ \le\ \TailBudget(X).
\]
\end{theorem}

\begin{proof}
By Theorem~\ref{thm:X.bankpositivity}, $\LB(X)\le\MainSym(X,n)$.
By Lemma~\ref{lem:X.noPProuting}, $\MainSym(X,n)\le \MainTail(X,n)+\ErrBudget(X,H,Q)$.
Hence $\Margin(X)=\LB-\ErrBudget\le \MainTail(X,n)$.
Apply Lemma~\ref{lem:T3-body} to bound $\MainTail(X,n)\le \TailBudget(X)$.
\end{proof}

\section{Machine check II: Section 11 endpoints and final assembly}
\addcontentsline{toc}{section}{Appendix Y. Section 11 endpoints and final assembly}

\subsection*{Endpoints}

\begin{theorem}[Type–I endpoint]\label{thm:typeIendpoint-body}
\(\ErrTypeI(X,H,Q)\le L^{C}\cdot X/H\) for \(X\ge X_0\).
\end{theorem}

\begin{proof}
\textbf{(TT\(^*\) mapping)} AO yields \(\PreErrTypeI\ll L^{C}\cdot E\) where \(E\) is the CLS energy.
\textbf{(Low-pass)} For \(\KH\) we have \(E\le (X+Q^2)/H\).
\textbf{(Ledger)} \(Q^2\le X\) gives \((X+Q^2)/H\le X/H\).
Thus \(\ErrTypeI\le \PreErrTypeI\ll L^{C}\cdot X/H\).
\end{proof}

\begin{theorem}[SSU endpoint]\label{thm:ssuendpoint-body}
\(\ErrSSU(X,H,Q)\le L^{-2}\) for \(X\ge X_0\) under \(\Aexp=10,\gexp=\tfrac{3}{10}\).
\end{theorem}

\begin{proof}
SSU + Schur bound: \(\ErrSSU\ll L^{C}\sqrt{H/X}\).
With \(H=(\log X)^{10}\) and \(X\ge 10^{12}\): \(\sqrt{H/X}\le 10^{-6}L^5\ll L^{-3}\).
Absorb \(L^C\) into the exponent (our \(A=10\) is ample) to get \(L^{-2}\).
\end{proof}

\begin{theorem}[Low-pass endpoint]\label{thm:lowpassendpoint-body}
\(\ErrTail(X,H,Q)\le \kappaFour/H^2\) for \(H\ge 1\).
\end{theorem}

\begin{proof}
Appendix A’s calculus for the \(\cos^2\) low-pass gives two integrations by parts and band-limit, hence \(O(H^{-2})\) with absolute \(\kappaFour\).
\end{proof}

\subsection*{Structural inequality and assembly}
\begin{theorem}[Structural decomposition]\label{thm:structural-body}
For even \(n\in[X,2X]\),
\(
R_2(X,n)\ \ge\ \MainSym(X,n)-(\ErrTypeI+\ErrSSU+\ErrTail).
\)
\end{theorem}

\begin{proof}
Bank/TFA algebra splits the correlation into the main and three error pieces; take absolute values on the error side.
\end{proof}

\begin{lemma}[Budget combiner]\label{lem:combinerY-body}
For \(X\ge X_0\),
\(\ErrBudget(X,H,Q)\le \frac{A_*+B_*+C_*}{L^2}\) with absolute \(A_*,B_*,C_*>0\).
\end{lemma}

\begin{proof}
Sum Theorems~\ref{thm:typeIendpoint-body}, \ref{thm:ssuendpoint-body}, \ref{thm:lowpassendpoint-body} under the ledger choice, absorbing constants.
\end{proof}

\begin{theorem}[Dyadic lower bound]\label{thm:dyadicLB-body}
For even \(n\in[X,2X]\), \(R_2(X,n)\ge \czero L^2-\frac{A_*+B_*+C_*}{L^2}\).
\end{theorem}

\begin{proof}
Section 10 positivity gives \(\MainSym\ge \LB=\czero L^2\). Apply Theorem~\ref{thm:structural-body} and Lemma~\ref{lem:combinerY-body}.
\end{proof}

\subsection*{Margin link}
Let \(\Margin(X)=\LB(X)-\ErrBudget(X,H,Q)\).
By Lemma~\ref{lem:combinerY-body}, \(\Margin(X)>0\) for \(X\ge X_0\).
Together with Theorem~\ref{thm:tailMarginBridge-body}\, this yields the dyadic-to-pointwise upgrade needed for the final certification.

\begin{lemma}[TT$^\ast$ mapping]\label{lem:Y.preErr.to.CLS}
There exists $C>0$ such that
\[
  \PreErrTypeI(X,H,Q)\ \le\ (\log(\max\{X,3\}))^{C}\cdot E_{\mathrm{CLS}}(X,H,Q),
\]
where $E_{\mathrm{CLS}}$ is the Gram–diagonalized TT$^\ast$ energy.
\end{lemma}

\begin{proof}
Write the Type–I bilinear form $B(f)$ and its TT$^\ast$ operator $T=B^\ast B$.
AO dispersion gives $\|G-I\|\le \varepsilon(H,Q)$ with $\varepsilon<1$, hence by Neumann series
$G$ is invertible and $\langle Gv,v\rangle\ge (1-\varepsilon)\|v\|^2$ for all $v$.
The occupancy bound (no hot spots) contributes $Q^{-2}$. Unfolding $T$ on bank atoms
and applying Cauchy–Schwarz yields
\(
  \PreErrTypeI \le (\log X)^{C}\langle Gv,v\rangle
   \le (\log X)^{C} E_{\mathrm{CLS}}.
\)
\end{proof}

\begin{lemma}[Low–pass squeeze]\label{lem:Y.KH.lowpass}
For the $\cos^2$ low–pass $K_H$,
\[
  E_{\mathrm{CLS}}(X,H,Q)\ \le\ \frac{X+Q^2}{H}.
\]
\end{lemma}

\begin{proof}
The Fourier support of $K_H$ lies in $|\xi|\le 1/H$ and $\int K_H=1$.
The CLS energy integrates the kernel along short shifts; restricting to $|\xi|\le 1/H$
forces a $1/H$ scaling. The $X$ term is the bulk, while the $Q^2$ term accounts for
aliasing up to denominator $Q$; both are standard in Appendix~A’s calculus.
\end{proof}

\begin{proof}[Proof of Theorem~\ref{thm:typeIendpoint-body}]
Combine Lemma~\ref{lem:Y.preErr.to.CLS} and Lemma~\ref{lem:Y.KH.lowpass}, then use the ledger constraint
$Q^2\le X$ to bound $(X+Q^2)/H\le X/H$.
\end{proof}

\begin{proof}[Proof of Theorem~\ref{thm:ssuendpoint-body}]
Decompose the off–diagonal into tubes $T_\tau$.
For each tube, Schur’s test with weights $w(d,n)=1$ gives
$\|T_\tau\|_{\ell^2\to\ell^2}\le \sqrt{(\sup_{(d,n)}\sum_{(d',n')}(K_H)_{\tau})\cdot
(\sup_{(d',n')}\sum_{(d,n)}(K_H)_{\tau})}$.
Low–pass smoothness and tube geometry produce $(\log X)^{C}\sqrt{H/X}$.
Summing over $\tau$ contributes another polylog factor, still absorbed into $(\log X)^C$.
With $H=(\log X)^{10}$ and $X\ge 10^{12}$, $\sqrt{H/X}\le 10^{-6} L^5\ll L^{-3}$, hence
$\ErrSSU\le L^{-2}$.
\end{proof}

\begin{lemma}[Mean–zero]\label{lem:Y.delta1}
The $\delta=1$ mask has zero mean: $\int_{\mathbb T}\!\delta_1(\theta)\,d\theta=0$.
\end{lemma}
\begin{proof}
By construction $\delta_1(\theta)=\psi(\theta)-\psi(\theta+1)$ with $\psi$ periodic and
$\int \psi= \int \psi(\cdot+1)$, hence their difference integrates to $0$.
\end{proof}

\begin{lemma}[Alias cancellation]\label{lem:Y.delta2}
The $\delta=2$ mask is row/column balanced; the first alias sum vanishes.
\end{lemma}
\begin{proof}
Arrange the bank atoms in a matrix by row/column indices; the balanced row/column sums
cancel the first harmonic by symmetry of the bank centers and the guard.
\end{proof}

\begin{lemma}[Lower frame from Gram closeness]\label{lem:Y.frame}
If $\|G-I\|\le \varepsilon<1$ for the bank Gram matrix, then
$\langle Gv,v\rangle\ge (1-\varepsilon)\|v\|^2$ for all $v$.
\end{lemma}
\begin{proof}
Neumann series: $G=I-(I-G)$ with $\|I-G\|<1$ implies $G$ invertible and
$\|G^{-1}\|\le 1/(1-\varepsilon)$. Then $(1-\varepsilon)\|v\|^2\le \langle Gv,v\rangle$.
\end{proof}

\begin{lemma}[Ledger monomial $\to$ log profile]\label{lem:Y.ledger}
For $X\ge X_0$, the ledger function $\Phi(X)$ satisfies
$\Phi(X)\le C_\Phi / (\log(\max\{X,3\}))^2$.
\end{lemma}
\begin{proof}
The concrete ledger inequalities reduce $\Phi$ to a finite maximum of rational expressions
in $X$ and $L=\log(\max\{X,3\})$. Each term is monotone in $X$ for $X\ge X_0$ and bounded by
$C_\Phi/L^2$; the constant $C_\Phi$ is fixed by the numeric ledger (Appendix tables).
\end{proof}
